%% ma: L12-B14-0010
%% lop: 12
%% chuong: 5
%% bai: 14
%% dang: 12.14.3
%% loai: TN4
%% muc: TH
%% dap_an: "A"
%% nguon: "(NBV)-ĐỀ SỐ 13-MỖI NGÀY 1 ĐỀ THI 2025.pdf (D13, MỖI NGÀY 1 ĐỀ - Đề số 13), Phần I Câu 12"
%% kien_thuc: "Bài 14 (vectơ pháp tuyến của mặt phẳng, tích có hướng)"
%% trang_thai: cho-duyet
%% ly_do: "Đề gốc có hai phương án đúng (A và C cùng phương), đã sửa phương án C, cần giáo viên xác nhận"
%% ghi_chu: "Đề gốc có hai phương án đều đúng: vectơ (3;-2;-1) ở A và vectơ (-3;2;1) ở C là hai vectơ cùng phương. Đã sửa phương án C thành (3;2;1) để chỉ còn đáp án A (đáp án gốc A). Cần giáo viên xác nhận cách sửa."
\begin{ex}
Trong không gian $Oxyz$, mặt phẳng đi qua $A(1;-3;5)$, $B(-1;-5;3)$ và $C(2;-1;4)$ có một vectơ pháp tuyến là
\choice{\True $\vec n=(3;-2;-1)$}{$\vec n=(3;2;-1)$}{$\vec n=(3;2;1)$}{$\vec n=(3;-2;-5)$}
\loigiai{Ta có $\overrightarrow{AB}=(-2;-2;-2)$, $\overrightarrow{AC}=(1;2;-1)$ và
\[\left[\overrightarrow{AB},\overrightarrow{AC}\right]=(6;-4;-2)=2(3;-2;-1).\]

Mặt phẳng $(ABC)$ có một vectơ pháp tuyến là $\vec n=(3;-2;-1)$. Chọn A.}
\end{ex}
